\section[]{ نرم ماتریس زیرگوسی}
فرض کنید $A$ یک ماتریس تصادفی
$m\times n$
با ردیف‌های 
$A_i$
باشد که $A_i$ ها بردارهای مستقل، با میانگین صفر، زیرگوسی و ایزوتروپیک هستند. نشان‌ دهید که برای هر $t$ مثبت، عبارت زیر با احتمال حداقل
$1 - 2 \exp(-t^2)$
برقرار است:
$$|| \frac{1}{m} A^TA - I_n||_{op} \leq C \max(\delta, \delta^2)$$
که در آن
$\delta = \sqrt{\frac{n}{m}} + \frac{t}{\sqrt{m}}$
است.

\vspace{0.5cm}

\textbf{یادداشت}:
بردار تصادفی $X$ زیرگوسی است اگر به ازای هر بردار $x$ عضو کره‌ی واحد، ضرب داخلی $X$ و $x$ زیرگوسی با ثابت محدود باشد. در قضیه‌ی فوق، ثابت‌های $C$ می‌توانند به ضرایب زیرگوسی بودن بردار $X$ نیز وابسته باشند.

\vspace{0.5cm}

\textbf{راهنمایی}:
نرم اپراتوری را به فرم سوپریمی آن نوشته و تلاش کنید با انتخاب یک 
\lr{$\epsilon$-net}
مناسب از کره‌ی واحد و با ایده‌ی گسسته‌سازی، کرانی برای احتمال مطلوب بیابید. از نامساوی برن‌اشتاین بهره ببرید. توجه کنید که تنها یک مرحله گسسته‌سازی کافی است و نیازی به استفاده از ایده‌های چینینگ نیست.


